% Source: https://artofproblemsolving.com/community/c3h2794491_2006__midmichigan_mathematical_olympiad___grades_56
% Stable identity: aops-topic-2794491-q03
The second player reflects each first-player placement across the vertical midline. No cell is fixed by this reflection. Reflected cells have opposite chessboard colours, so the paired bishops cannot attack each other. Before every first-player turn, the existing configuration is reflection-symmetric. If a new legal cell is unattacked and empty, its reflected cell also is; placing its partner preserves symmetry. Thus the second player always has a reply and wins when the finite board leaves the first with no move.
\begin{asy}
usepackage("amsmath");
size(7cm);
for(int k=0;k<=8;++k){draw((k,0)--(k,8),gray);draw((0,k)--(8,k),gray);}draw((4,0)--(4,8),dashed);dot((1.5,2.5),blue);dot((6.5,2.5),blue);label("$P$",(1.5,2.5),N);label("$P^\prime$",(6.5,2.5),N);
\end{asy}
